% Part: computability
% Chapter: computability-theory
% Section: equiv-ce-defs

\documentclass[../../../include/open-logic-section]{subfiles}

\begin{document}

\olfileid{cmp}{thy}{eqc}

\olsection[પરિગણનીય ગણોની વ્યાખ્યાઓ]{સંગણકીય રીતે પરિગણનીય
  ગણોની સમકક્ષ વ્યાખ્યાઓ}


આગળની વાત સંગણકીય રીતે પરિગણનીય હોવાના અર્થનાં કેટલાંક
મહત્ત્વનાં સમકક્ષ વિધાનો આપે છે.

\begin{thm}
\ollabel{thm:ce-equiv}
માનો કે $S$ પ્રાકૃતિક સંખ્યાઓનો ગણ છે. તો નીચેના વિધાનો
સમકક્ષ છે:
\begin{enumerate}
\item\ollabel{case:ce} $S$ સંગણકીય રીતે પરિગણનીય છે.
\item\ollabel{case:ran-pc} $S$ કોઈ \emph{આંશિક} સંગણનીય વિધેયનો મૂલ્યગણ છે.
\item\ollabel{case:ran-prim} $S$ ખાલી છે અથવા આદિમ પુનરાવર્તી વિધેયનો મૂલ્યગણ છે.
\item\ollabel{case:ce-domain} $S$ કોઈ આંશિક સંગણનીય વિધેયનો \emph{પ્રદેશ} છે.
\end{enumerate}
\end{thm}

\begin{explain}
પહેલાં ત્રણ ખંડો કહે છે કે દરેક ખાલી ન હોય એવા સંગણકીય રીતે
પરિગણનીય ગણની પરિગણના સંગણનીય વિધેય, આંશિક સંગણનીય વિધેય
અથવા આદિમ પુનરાવર્તી વિધેય વડે થઈ શકે છે. ચોથો ખંડ કહે છે કે
$S$ સંગણકીય રીતે પરિગણનીય હોય, તો કોઈ સૂચક~$e$ માટે
\[
S = \Setabs{x}{\cfind{e}(x) \fdefined}.
\]
બીજા શબ્દોમાં, $S$ એ એવા ઇનપુટોનો ગણ છે જેના માટે
$\cfind{e}$ની સંગણના થંભે છે. તેથી સંગણકીય રીતે પરિગણનીય
ગણોને ક્યારેક \emph{અર્ધનિર્ણેય} પણ કહે છે: સંખ્યા ગણમાં હોય
તો છેવટે ``હા'' મળે છે, પરંતુ ન હોય તો ``ના'' કદી મળતું નથી!{}
\end{explain}

\begin{proof}
દરેક આદિમ પુનરાવર્તી વિધેય સંગણનીય છે અને દરેક સંગણનીય વિધેય
આંશિક સંગણનીય પણ છે. તેથી \olref{case:ran-prim} પરથી
\olref{case:ce} અને \olref{case:ce} પરથી
\olref{case:ran-pc} મળે છે. (જો $S$ ખાલી હોય, તો $S$ ક્યાંય
વ્યાખ્યાયિત ન હોય એવા આંશિક સંગણનીય વિધેયનો મૂલ્યગણ છે.)
\olref{case:ran-pc} પરથી \olref{case:ran-prim} બતાવીએ,
તો પહેલાં ત્રણ ખંડોની સમકક્ષતા સાબિત થશે.

હવે માનો કે $S$ આંશિક સંગણનીય વિધેય $\cfind{e}$નો મૂલ્યગણ
છે. $S$ ખાલી હોય તો કામ પૂરું. અન્યથા $S$નું કોઈ તત્ત્વ $a$
લો. ક્લીનીના સામાન્ય સ્વરૂપના પ્રમેયથી
\[
\cfind{e}(x) = U(\umin{s}{T(e, x, s)}).
\]
ખાસ કરીને, $\cfind{e}(x) \fdefined$ અને તેનું મૂલ્ય $= y$
ત્યારે અને માત્ર ત્યારે જ્યારે કોઈ $s$ માટે $T(e, x, s)$
અને $U(s) = y$ હોય. $f(z)$ને નીચે પ્રમાણે વ્યાખ્યાયિત કરો:
\[
f(z) = \begin{cases}
  U((z)_1) & \text{જો $T(e, (z)_0, (z)_1)$} \\
  a        & \text{અન્યથા.}
\end{cases}
\]
$T$ અને $U$ આદિમ પુનરાવર્તી હોવાથી $f$ પણ આદિમ પુનરાવર્તી
છે. \iftag{TMs}{ટ્યુરિંગ મશીનની ભાષામાં, જો $z$ એવી
  જોડી $\tuple{(z)_0, (z)_1}$નો કોડ હોય કે $(z)_1$,
  મશીન~$M_e$ની ઇનપુટ $(z)_0$ પરની થંભતી સંગણના હોય,
  તો $f$ તે સંગણનાનું આઉટપુટ આપે છે; અન્યથા~$a$ આપે છે.}

હવે $S$, $f$નો મૂલ્યગણ છે તે બતાવવાનું છે: દરેક પ્રાકૃતિક
સંખ્યા~$y$ માટે $y \in S$ ત્યારે અને માત્ર ત્યારે જ્યારે
તે $f$ના મૂલ્યગણમાં હોય. એક દિશામાં, માનો કે $y \in S$.
તો $y$, $\cfind{e}$ના મૂલ્યગણમાં છે; તેથી કોઈ $x$ અને~$s$
માટે $T(e,x,s)$ અને $U(s) = y$ છે. પરંતુ પછી
$y = f(\tuple{x,s})$. ઊલટી દિશામાં, માનો કે $y$, $f$ના
મૂલ્યગણમાં છે. તો $y = a$ છે અથવા કોઈ~$z$ માટે
$T(e,(z)_0,(z)_1)$ અને $U((z)_1) = y$ છે. પછીના કિસ્સામાં
$\cfind{e}((z)_0) \fdefined = y$ હોવાથી બંને રીતે $y$,
$S$માં છે.
\sourcecorrection{OLCT-004}{સ્થિર અંગ્રેજી મૂળ છેલ્લી
અનુમાન-કડીમાં $\cfind{e}(x) \fdefined = y$ લખે છે,
પરંતુ આ દિશામાં ઉપલબ્ધ સાક્ષી $z$ છે; તેનું પ્રથમ ઘટક
$(z)_0$ ઇનપુટ છે. તેથી અહીં $\cfind{e}((z)_0)$ રાખ્યું છે.}

(સંકેત $\cfind{e}(x) \fdefined = y$નો અર્થ
``$\cfind{e}(x)$ વ્યાખ્યાયિત છે અને તેનું મૂલ્ય $y$ છે.''
અહીં $\cfind{e}(x) = y$ પણ લખી શકાય; વધારાનું તીર
આંશિક વિધેય સાથે કામ કરીએ છીએ એ યાદ અપાવે છે.)

\olref{thm:ce-equiv}ની સાબિતી પૂરી કરવા
\olref{case:ce} અને~\olref{case:ce-domain} સમકક્ષ છે એ
બતાવવું પૂરતું છે. પહેલાં \olref{case:ce} પરથી
\olref{case:ce-domain} બતાવીએ. ખાલી ગણ ક્યાંય વ્યાખ્યાયિત
ન હોય એવા આંશિક સંગણનીય વિધેયનો પ્રદેશ છે; હવે બાકીના
કિસ્સામાં માનો કે $S$ સંગણનીય વિધેય~$f$નો મૂલ્યગણ છે, એટલે
\[
S = \Setabs{y}{\text{કોઈ $x$ માટે, } f(x) = y}.
\]
\sourcecorrection{OLCT-005}{સ્થિર અંગ્રેજી સાબિતી આ
દિશામાં ખાલી ગણને અલગથી લેતી નથી, જ્યારે કુલ વ્યાખ્યા
તેનો સમાવેશ કરે છે. અહીં ખાલી ગણ માટે ક્યાંય વ્યાખ્યાયિત
ન હોય એવા આંશિક વિધેયનો પ્રદેશ સ્પષ્ટ કર્યો છે.}
માનો કે
\[
g(y) = \umin{x}{(f(x) = y)}.
\]
ત્યારે $g$ આંશિક સંગણનીય છે અને $g(y)$ ત્યારે અને માત્ર
ત્યારે વ્યાખ્યાયિત છે જ્યારે કોઈ~$x$ માટે $f(x) = y$.
બીજા શબ્દોમાં, $g$નો પ્રદેશ $f$નો મૂલ્યગણ છે.
\iftag{TMs}{ટ્યુરિંગ મશીનની ભાષામાં: $S$નાં તત્ત્વોની
પરિગણના કરતું ટ્યુરિંગ મશીન~$F$ હોય, તો $G$ એવું મશીન
બનાવો કે જે $F$નાં આઉટપુટોમાં આપેલું તત્ત્વ શોધીને $S$નું
અર્ધનિર્ણય કરે; તત્ત્વ મળે તો થંભે અને ન મળે તો અનંતકાળ સુધી
શોધ ચાલુ રાખે.}{}

છેવટે \olref{case:ce-domain} પરથી~\olref{case:ce}
બતાવવા, માનો કે $S$ આંશિક સંગણનીય વિધેય~$\cfind{e}$નો
પ્રદેશ છે, એટલે
\[
S = \Setabs{x}{\cfind{e}(x) \fdefined}.
\]
$S$ ખાલી હોય તો કામ પૂરું; અન્યથા $S$નું કોઈ તત્ત્વ $a$
લો. $f$ને નીચે પ્રમાણે વ્યાખ્યાયિત કરો:
\[
f(z) = \begin{cases}
(z)_0 & \text{જો $T(e,(z)_0,(z)_1)$} \\
a & \text{અન્યથા.}
\end{cases}
\]
તો, ઉપરની જેમ, સંખ્યા $x$, $f$ના મૂલ્યગણમાં ત્યારે અને
માત્ર ત્યારે છે જ્યારે $\cfind{e}(x) \fdefined$, એટલે જ્યારે
$x \in S$. \iftag{TMs}{ટ્યુરિંગ મશીનની ભાષામાં:
$S$નું અર્ધનિર્ણય કરતું મશીન $M_e$ હોય, તો ટ્યુરિંગ મશીનની
બધી શક્ય સંગણનાઓ તપાસીને, થંભતી સંગણનાઓને અનુરૂપ
ઇનપુટો આપતાં જવાથી $S$નાં તત્ત્વોની પરિગણના થાય છે.}{}
\end{proof}

\olref{thm:ce-equiv}ના ખંડ~\olref{case:ce-domain} વડે
સંગણકીય રીતે પરિગણનીય ગણોની પરિગણના કરવાની સગવડભરી રીત
મળે છે: દરેક~$e$ માટે $W_e$, $\cfind{e}$નો પ્રદેશ હોય, એટલે
\[
W_e = \Setabs{x}{\cfind{e}(x) \fdefined}.
\]
તો $A$ કોઈ પણ સંગણકીય રીતે પરિગણનીય ગણ હોય, તો કોઈ~$e$
માટે $A = W_e$.

આગળનું પ્રમેય સંગણકીય રીતે પરિગણનીય ગણોનું વધુ એક લક્ષણ આપે છે.

\begin{thm}
\ollabel{thm:exists-char}
ગણ $S$ ત્યારે અને માત્ર ત્યારે સંગણકીય રીતે પરિગણનીય છે
જ્યારે એવો સંગણનીય સંબંધ $R(x,y)$ હોય કે
\[
S = \Setabs{ x }{ \lexists[y][R(x,y)] }.
\]
\end{thm}

\begin{proof}
એક દિશામાં, માનો કે $S$ સંગણકીય રીતે પરિગણનીય છે.
તો કોઈ $e$ માટે $S = W_e$. આ~$e$ માટે $S$ને આમ લખી શકાય:
\[
S = \Setabs{ x }{ \lexists[y][T(e, x, y)] }.
\]
ઊલટી દિશામાં, માનો કે $S = \Setabs{ x }{
  \lexists[y][R(x, y)] }$. $f$ને આ રીતે વ્યાખ્યાયિત કરો:
\[
f(x) \simeq \umin{y}{R(x, y)}.
\]
તો $f$ આંશિક સંગણનીય છે અને $S$, $f$નો પ્રદેશ છે.
\end{proof}


\end{document}
